\clearpage
\section{模型假设}

为使问题在数学上可处理，同时保留对实际情况的合理描述，本文提出以下基本假设。

\textbf{假设1（理想平面纸张）：} 纸张为严格平面，位于 $z=0$ 平面，无弯曲褶皱。题目明确给定"理想平面"，A4纸在平放状态下可视为平面，该假设与题目条件一致。实际中A4纸的翘曲量级约为0.1--0.5 mm，相对于圆柱半径（10--30 mm）可忽略不计，因此该假设引入的误差极小。

\textbf{假设2（理想圆柱镜面）：} 镜面为完美圆柱体侧面，表面光滑，反射率均匀。题目明确给定"理想圆柱体侧面"，忽略实际镜面的微小瑕疵和边缘效应，使得反射定律可精确应用于圆柱面上的每一个点。工业级镜面圆柱的表面粗糙度通常小于0.1 $\mu$m，远小于可见光波长，因此镜面反射假设成立。

\textbf{假设3（几何光学近似）：} 光线传播遵循几何光学定律，忽略衍射和散射效应。可见光波长（400--700 nm）远小于圆柱尺寸（厘米量级），菲涅尔数 $F = a^2/(\lambda L) \gg 1$（其中 $a$ 为圆柱半径，$L$ 为观察距离），几何光学近似的误差可忽略不计。该假设使得反射问题可以用射线追踪方法精确求解。

\textbf{假设4（远场平行光近似）：} 观察者距离足够远，入射光线近似为平行光。实际观察距离（约0.5--1 m）远大于圆柱半径（约1--3 cm），平行光近似引入的角度误差约为 $\arctan(R/L) \approx R/L \approx 0.02$ rad，对应的纸面位置偏差小于1 mm，在工程精度范围内可接受。该假设的核心优势在于使得所有入射光线共享同一方向向量 $\boldsymbol{v}$，大幅简化了逆映射公式的推导。

\textbf{假设5（圆柱竖直放置）：} 圆柱轴线严格垂直于纸面，底面与纸面接触。题目给定"竖直放置"，该假设简化了三维几何关系，使逆映射公式具有闭式解。若圆柱轴线倾斜，法向量将包含 $z$ 分量，反射公式的复杂度将显著增加，但基本建模框架仍然适用。

\textbf{假设6（标准双目观察者）：} 观察者为正常视力的成年人，瞳距 $d_{\text{IPD}} = 64$ mm（成人平均值），人眼等效焦距 $f_e = 17$ mm，从斜上方以双眼观察圆柱镜面，观察距离约为 $400\sim800$ mm。该假设引入了人因约束，使模型能够分析双目视差、舒适视区和复视阈值等实际观察体验因素。
